
\documentclass{article}
\usepackage{colortbl}
\usepackage{makecell}
\usepackage{multirow}
\usepackage{supertabular}

\begin{document}

\newcounter{utterance}

\centering \large Interaction Transcript for game `referencegame', experiment `random\_grids', episode 8 with reflexion{-}llama{-}3.1{-}8b{-}instruct.
\vspace{24pt}

{ \footnotesize  \setcounter{utterance}{1}
\setlength{\tabcolsep}{0pt}
\begin{supertabular}{c@{$\;$}|p{.15\linewidth}@{}p{.15\linewidth}p{.15\linewidth}p{.15\linewidth}p{.15\linewidth}p{.15\linewidth}}
    \# & $\;$A & \multicolumn{4}{c}{Game Master} & $\;\:$B\\
    \hline

    \theutterance \stepcounter{utterance}  
    & & \multicolumn{4}{p{0.6\linewidth}}{
        \cellcolor[rgb]{0.9,0.9,0.9}{
            \makecell[{{p{\linewidth}}}]{
                \texttt{\tiny{[P1$\langle$GM]}}
                \texttt{You have attempted this same task before and did not succeed. These are the plans you wrote after those attempts.} \\
\texttt{Trial \#1: New plan: I failed to distinguish the target grid from the distractors because I focused on the rows with two X's, which are present in both distractor grids as well. Next time, I will look for a unique feature that is present in the target grid but not in the distractors. I will examine the columns and look for a combination of X's and empty cells that is unique to the target grid. For example, I will check if there is a column with a specific pattern of X's and empty cells that is not present in the distractors.} \\
\\ 
\texttt{Trial \#2: New plan: I failed to distinguish the target grid from the distractors because I focused on a column with a specific pattern of X's and empty cells, but I didn't consider the possibility that this pattern might be present in the distractors as well. Next time, I will look for a unique feature that is present in the target grid but not in the distractors. I will examine the rows and look for a combination of X's and empty cells that is unique to the target grid. For example, I will check if there is a row with a specific pattern of X's and empty cells that is not present in the distractors. Additionally, I will also consider the placement of X's in relation to the edges of the grid, as this might provide a unique feature that distinguishes the target grid from the distractors.} \\
\\ 
\\ 
\texttt{You are given three grids, where each of them is 5 by 5 in size.} \\
\texttt{Grids have empty cells marked with "▢" and filled cells marked with "X".} \\
\texttt{Your task is to generate a referring expression that best describes the target grid while distinguishing it from the two other distractor grids.} \\
\texttt{The first grid is the target grid, and the following two grids are the distractors.} \\
\\ 
\texttt{Target grid:} \\
\\ 
\texttt{□ □ X X □} \\
\texttt{X □ □ X X} \\
\texttt{□ X X □ X} \\
\texttt{X X □ □ X} \\
\texttt{□ X X □ □} \\
\\ 
\texttt{Distractor grid 1:} \\
\\ 
\texttt{X □ X X □} \\
\texttt{□ X □ X X} \\
\texttt{X X □ □ X} \\
\texttt{□ X X □ X} \\
\texttt{X □ X □ □} \\
\\ 
\texttt{Distractor grid 2:} \\
\\ 
\texttt{□ X X □ X} \\
\texttt{X □ X X □} \\
\texttt{□ X □ X X} \\
\texttt{X □ X □ X} \\
\texttt{□ X X □ □} \\
\\ 
\texttt{Instruction: Describe the target grid.} \\
\texttt{Generate the referring expression starting with the tag "Expression: " for the given target grid. Omit any other text.} \\
            }
        }
    }
    & \\ \\

    \theutterance \stepcounter{utterance}  
    & \multicolumn{4}{p{0.6\linewidth}}{
        \cellcolor[rgb]{0.9,0.9,0.9}{
            \makecell[{{p{\linewidth}}}]{
                \texttt{\tiny{[P1$\rangle$GM]}}
                \texttt{Expression: The grid that has a row with three X's in the top row and a row with two X's in the third row from the top, and a column with an X in the first row and an X in the fourth row.} \\
            }
        }
    }
    & & \\ \\

    \theutterance \stepcounter{utterance}  
    & & & \multicolumn{2}{p{0.3\linewidth}}{
        \cellcolor[rgb]{0.9,0.9,0.9}{
            \makecell[{{p{\linewidth}}}]{
                \texttt{\tiny{[GM$|$GM]}}
                \texttt{Expression: The grid that has a row with three X's in the top row and a row with two X's in the third row from the top, and a column with an X in the first row and an X in the fourth row.} \\
            }
        }
    }
    & & \\ \\

    \theutterance \stepcounter{utterance}  
    & & \multicolumn{4}{p{0.6\linewidth}}{
        \cellcolor[rgb]{0.9,0.9,0.9}{
            \makecell[{{p{\linewidth}}}]{
                \texttt{\tiny{[GM$\rangle$P2]}}
                \texttt{You have attempted this same task before and did not succeed. These are the plans you wrote after those attempts.} \\
\texttt{Trial \#1: New plan: I failed to identify the correct grid because I didn't carefully analyze the expression. The expression mentions a row with two X's in the first and last columns, but I didn't notice that the first and last columns are not the same in all grids. I also didn't notice that the second and fourth columns are not the same in all grids. Next time, I will carefully read the expression and look for specific features that distinguish the correct grid from the others. I will also pay attention to the fact that the first and last columns, as well as the second and fourth columns, are not necessarily the same in all grids.} \\
\\ 
\texttt{Trial \#2: New plan: I failed to identify the correct grid because I didn't carefully analyze the expression. The expression mentions a column with an X in the first row, an X in the third row, and an empty cell in the fifth row. However, I didn't notice that the column with these features is not unique to one grid. In fact, the first grid has this column, but the second grid also has a column with an X in the first row, an X in the third row, and an empty cell in the fifth row. Next time, I will carefully read the expression and look for specific features that distinguish the correct grid from the others. I will also pay attention to the fact that the column with the specified features might not be unique to one grid. I will look for other features that are mentioned in the expression, such as the column with an X in the first row, and use them to narrow down the possible grids.} \\
\\ 
\\ 
\texttt{You are given three grids, where each of them is 5 by 5 in size.} \\
\texttt{Grids have empty cells marked with "▢" and filled cells marked with "X".} \\
\texttt{You are also given a referring expression that describes one of the given grids.} \\
\texttt{Your task is to select the grid that matches the given referring expression.} \\
\texttt{Generate only the number (in text) of the grid that matches the given expression by selecting first, second, or third.} \\
\\ 
\texttt{First:} \\
\\ 
\texttt{□ X X □ X} \\
\texttt{X □ X X □} \\
\texttt{□ X □ X X} \\
\texttt{X □ X □ X} \\
\texttt{□ X X □ □} \\
\\ 
\texttt{Second:} \\
\\ 
\texttt{X □ X X □} \\
\texttt{□ X □ X X} \\
\texttt{X X □ □ X} \\
\texttt{□ X X □ X} \\
\texttt{X □ X □ □} \\
\\ 
\texttt{Third:} \\
\\ 
\texttt{□ □ X X □} \\
\texttt{X □ □ X X} \\
\texttt{□ X X □ X} \\
\texttt{X X □ □ X} \\
\texttt{□ X X □ □} \\
\\ 
\texttt{Expression: The grid that has a row with three X's in the top row and a row with two X's in the third row from the top, and a column with an X in the first row and an X in the fourth row.} \\
\texttt{Question: Which grid does the expression refer to?} \\
\texttt{Start with the tag "Answer: ", followed by your selection. Omit any other text.} \\
            }
        }
    }
    & \\ \\

    \theutterance \stepcounter{utterance}  
    & & & \multicolumn{4}{p{0.6\linewidth}}{
        \cellcolor[rgb]{0.9,0.9,0.9}{
            \makecell[{{p{\linewidth}}}]{
                \texttt{\tiny{[GM$\langle$P2]}}
                \texttt{Answer: first} \\
            }
        }
    }
     \\ \\

    \theutterance \stepcounter{utterance}  
    & & & \multicolumn{2}{p{0.3\linewidth}}{
        \cellcolor[rgb]{0.9,0.9,0.9}{
            \makecell[{{p{\linewidth}}}]{
                \texttt{\tiny{[GM$|$GM]}}
                \texttt{first} \\
            }
        }
    }
    & & \\ \\

    \theutterance \stepcounter{utterance}  
    & \multicolumn{4}{p{0.6\linewidth}}{
        \cellcolor[rgb]{0.95,0.95,0.95}{
            \makecell[{{p{\linewidth}}}]{
                \texttt{\tiny{[P2$\rightarrow$P1]}}
                \texttt{Post{-}episode Reflexion (not sent to GM):} \\
\texttt{New plan: I failed to distinguish the target grid from the distractors because I focused on specific patterns of X's and empty cells in rows and columns, but I didn't consider the possibility that these patterns might be present in the distractors as well. Next time, I will look for a unique feature that is present in the target grid but not in the distractors. I will examine the target grid more carefully and look for a combination of X's and empty cells that is unique to the target grid. Additionally, I will also consider the placement of X's in relation to the edges of the grid, as this might provide a unique feature that distinguishes the target grid from the distractors. Furthermore, I will also consider the overall structure of the grid, such as the number of X's in each row and column, and look for a unique combination that is not present in the distractors.} \\
            }
        }
    }
    & & \\ \\

    \theutterance \stepcounter{utterance}  
    & & & \multicolumn{4}{p{0.6\linewidth}}{
        \cellcolor[rgb]{0.95,0.95,0.95}{
            \makecell[{{p{\linewidth}}}]{
                \texttt{\tiny{[P1$\rightarrow$P2]}}
                \texttt{Post{-}episode Reflexion (not sent to GM):} \\
\texttt{New plan: I failed to identify the correct grid because I didn't carefully analyze the expression. The expression mentions a row with three X's in the top row and a row with two X's in the third row from the top, and a column with an X in the first row and an X in the fourth row. However, I didn't notice that the column with an X in the first row and an X in the fourth row is not unique to one grid. In fact, the first and third grids both have this column. Next time, I will carefully read the expression and look for specific features that distinguish the correct grid from the others. I will also pay attention to the fact that the column with the specified features might not be unique to one grid. I will look for other features that are mentioned in the expression, such as the row with three X's in the top row and the row with two X's in the third row from the top, and use them to narrow down the possible grids. Specifically, I will look for the row with three X's in the top row and check which grid has this row. Then, I will look for the row with two X's in the third row from the top and check which grid has this row. By combining these two features, I should be able to identify the correct grid.} \\
            }
        }
    }
     \\ \\

\end{supertabular}
}

\end{document}
